\chapter{Direct sums and spectral Hall operations}
\label{cyspec:chapter}

Two operations on representations distinguish the geometry needed
for a spectral coproduct. The Hall correspondence remembers a
subrepresentation and integrates over its possible positions.
The direct-sum inclusion fixes two summands and sets both
off-diagonal arrow blocks to zero. Their normal bundles differ.

For one loop and two one-dimensional summands, write $x,y$ for
their Chern roots. In the Hall integral, the arrow normal weight
$y-x$ cancels the tangent weight of the invariant-line flag.
The two coordinate lines therefore give
$\mathbf 1_1\star\mathbf 1_1=2\mathbf 1_2$.
For the direct-sum inclusion, both off-diagonal entries vanish.
Their weights are $y-x$ and $x-y$, so its Gysin map sends
$1\otimes1$ to $-(y-x)^2$.

Here $\mathbf 1_n$ denotes the constant class in charge $n$.
Only $\mathbf 1_0$ is the Hall unit.

The direct-sum Gysin map becomes invertible away from its normal
Euler divisor. The Hall coproduct instead weights direct-sum
pullback by the Euler class of a quiver Hom complex. Translating
one summand makes its moving factors invertible as Laurent series.
Compatibility with the Hall product then involves four blocks:
each input has a part in each output.
We first construct these maps on polynomial coefficients, then
give their critical-cohomology form. An explicit pair of shuffle
algebras will show which collision information their common
exchange ratio fails to determine.

\section{Direct-sum normal bundles}
\label{cyspec:sec:direct-sums}

Work over $\mathbb C$ with rational cohomology. Let $Q$ be a finite
quiver, with vertex set $Q_0$ and arrow set $Q_1$.
A dimension vector $d=(d_i)_{i\in Q_0}$ has nonnegative integral entries.
Its representation space is
\[
 V_d=\bigoplus_{a:i\to j}\operatorname{Hom}(\mathbb C^{d_i},
                                           \mathbb C^{d_j}).
\]

Retain labeled roots and put
\[
 P_d=\mathbb Q[x_{i,r}:i\in Q_0,\ 1\leq r\leq d_i],
 \qquad C_d=\operatorname{Spec}P_d,
 \qquad m(d)=\sum_{a:i\to j}d_i d_j.
\]
Thus $P_d$ is the equivariant cohomology of $V_d$ for its coordinate
torus. A root has degree two. The space $C_d$ is the spectrum of the underlying
ordinary ring, with this grading retained on its coefficient modules.

For chosen blocks of dimensions $d,e$, the inclusion
$V_d\times V_e\hookrightarrow V_{d+e}$ has normal representation
\[
 N_{d,e}=\bigoplus_{a:i\to j}
 \left(\operatorname{Hom}(\mathbb C^{d_i},\mathbb C^{e_j})
 \oplus\operatorname{Hom}(\mathbb C^{e_i},\mathbb C^{d_j})\right).
\]
Consequently its Euler class is
\begin{equation}\label{cyspec:eq:direct-euler}
 \mathcal E_{d,e}(x;y)=
 \prod_{a:i\to j}
 \left(\prod_{r\leq d_i,\ s\leq e_j}(y_{j,s}-x_{i,r})\right)
 \left(\prod_{r\leq e_i,\ s\leq d_j}(x_{j,s}-y_{i,r})\right).
\end{equation}
Let $\mathcal F_d$ be the free rank-one graded sheaf on $C_d$,
where a polynomial $F$ has adjusted degree $\deg F-2m(d)$.
Under $C_{d+e}=C_d\times C_e$ for the chosen blocks, define
\begin{equation}\label{cyspec:eq:direct-map}
 \nu_{d,e}(F\boxtimes G)=\mathcal E_{d,e}(x;y)F(x)G(y).
\end{equation}

\begin{proposition}\label{cyspec:prop:direct-sums}
The maps $\nu_{d,e}$ have degree zero, are associative, and commute
with block exchange and root relabeling. Charge zero gives their unit.
On the open set where all cross-block roots are distinct, they are
isomorphisms. For any ordered collection of blocks, every iterated
map is
\begin{equation}\label{cyspec:eq:all-blocks}
 \boxtimes_{r=1}^k F_r\longmapsto
 \left(\prod_{r<s}\mathcal E_{d_r,d_s}(x_r;x_s)\right)
 \prod_{r=1}^k F_r(x_r).
\end{equation}
\end{proposition}

\begin{proof}
Complex orientation identifies the Gysin map with multiplication
by the Euler class of $N_{d,e}$. Its degree is
$2(m(d+e)-m(d)-m(e))$, which the three grading adjustments cancel.
For three blocks the normal representation is the direct sum of
the representations between each pair of blocks. Hence
\[
 \mathcal E_{d,e}\mathcal E_{d+e,f}
 =\mathcal E_{d,e}\mathcal E_{d,f}\mathcal E_{e,f}
 =\mathcal E_{e,f}\mathcal E_{d,e+f}.
\]

The same decomposition proves \eqref{cyspec:eq:all-blocks} for every
$k$. Both directions of each arrow occur, so exchanging the blocks
preserves the normal representation and its Euler class.
Relabeling roots permutes the factors. On the indicated open set
each factor is invertible, and division gives the inverse map.
An empty block contributes the empty product, equal to one.
\end{proof}

This is the direct-sum factorization over Chern-root parameters
\cite{cyspec:KS}*{\S6.5}. Its extension across the Euler divisor
contains more information than its restriction to that open set.

For one vertex with $\ell\geq1$ loops, take two charge-one blocks
and put $P=\mathbb Q[x,y]$, $\delta=y-x$. The Euler class is
$(-1)^\ell\delta^{2\ell}$. Forgetting the internal polynomial grading,
the cone of the normal map has the free representative
\begin{equation}\label{cyspec:eq:direct-cone}
 [\,P\xrightarrow{(-1)^\ell\delta^{2\ell}}P\,],
 \qquad\text{in cochain degrees }-1,0.
\end{equation}
It resolves $P/(\delta^{2\ell})$, since $P$ is a domain.

For $D=P/(\delta)$ its derived restriction is
$[D\xrightarrow{0}D]$. Thus the diagonal fiber has a class in each
of degrees $-1$ and $0$. Ordinary restriction of the degree-zero
quotient alone would omit the first class.

\section{The completion of an Euler inverse}
\label{cyspec:sec:euler-completion}

An external torus weight does not by itself invert the Euler class
of a bundle over a stack. Let $L$ be the universal line on
$B\mathbb G_m$, with $a=c_1(L)$, and let another torus $T=\mathbb G_m$
scale its fibers. Write $u$ for the external parameter.
The zero section has equivariant Euler class $a+u$.
The Gysin sequence for its complement $L^\times$ gives
\begin{equation}\label{cyspec:eq:universal-line}
 H_T^\bullet(L^\times;\mathbb Q)=\mathbb Q[a,u]/(a+u).
\end{equation}
There is no odd term because multiplication by $a+u$ is injective.

After tensoring over $\mathbb Q[u]$ with $\mathbb Q(u)$, the quotient
is still $\mathbb Q(u)$: the relation sets $a=-u$.
Equivalently, $u+a$ is not a unit of $\mathbb Q(u)[a]$.

The stabilizer accounts for the surviving class. Present the
complement as $[\mathbb G_m/H]$, where $H=\mathbb G_m$ acts by
scaling. In the further quotient by $T$, the point represented
by $1$ has automorphisms $(h,t)$ satisfying $ht=1$.
Their projection onto $T$ is an isomorphism. The free action on
each geometric fiber therefore does not remove the torus
stabilizer of the stack point.

For a commutative $\mathbb Q$-algebra $R$, write $R((u^{-1}))$ for series
whose powers of $u$ are bounded above, completed toward decreasing
powers. Coefficients may have unbounded degree in $R$.

\begin{proposition}\label{cyspec:prop:euler-inverse}
For $R=\mathbb Q[a]$, the Euler class $u+a$ is a unit of
$R((u^{-1}))$, with
\begin{equation}\label{cyspec:eq:euler-inverse}
 (u+a)^{-1}=u^{-1}\sum_{j\geq0}(-a/u)^j.
\end{equation}
On the finite approximation $\mathbb P^m$ of $B\mathbb G_m$, the
series stops at $j=m$. More generally, every factor $wu+c$, with
$0\ne w\in\mathbb Q$ and $c\in R$, is invertible in this completion.
\end{proposition}

\begin{proof}
Multiplication of the partial sum through $j=m$ by $u+a$ gives
$1-(-a/u)^{m+1}$. The remainder tends to zero in the specified
topology. On $\mathbb P^m$ it is zero because $a^{m+1}=0$.
Factoring $wu+c=wu(1+c/(wu))$ proves the last assertion.
\end{proof}

The corresponding flag denominator is equally concrete.
Let $T_2=\mathbb G_m^2$ act on $\mathbb C^2$, put
$R=\mathbb Q[x_1,x_2]$, and let $h$ be the first Chern class of
the tautological line on $\mathbb P^1$. Then
\[
 H_{T_2}^\bullet(\mathbb P^1)=R[h]/((h-x_1)(h-x_2)).
\]

Restriction to the two coordinate lines is
\[
 r(f_0+f_1h)=(f_0+x_1f_1,f_0+x_2f_1).
\]
Its matrix in the basis $1,h$ has determinant $\delta=x_2-x_1$.
Its image consists precisely of pairs $(a,b)$ with $b-a\in\delta R$.
After inverting $\delta$, its inverse is
\begin{equation}\label{cyspec:eq:flag-inverse}
 r^{-1}(a,b)=\frac{x_2a-x_1b}{\delta}+\frac{b-a}{\delta}h.
\end{equation}
Inverting $x_2$ alone cannot give this inverse: at $x_1=x_2=1$
the matrix has rank one.

For a bounded complex of bundles with nonzero external weights,
its Euler factor is a finite product of factors $wu+c$ and their
inverses. Proposition~\ref{cyspec:prop:euler-inverse} constructs this
characteristic class in the Laurent ring, independently of a
splitting into Chern roots. A Gysin inverse requires a further
identification: on the chosen fundamental class, the Gysin map
must act by this Euler factor.
Even coefficient extraction depends on the expansion: $u/(u+a)$
has constant term one at infinity and zero at the origin over
$\mathbb Q(a)$.

\section{The translated Hall coproduct}
\label{cyspec:sec:polynomial-coproduct}

Allow a torus $T$ to act on the arrows of $Q$, and let
$R=H^\bullet(BT;\mathbb Q)$. For an arrow $a$, denote its weight
class by $h_a\in R^2$. A zero weight is allowed.
For each dimension vector define the even polynomial coefficient
module
\[
 A_d=R[x_{i,r}:1\leq r\leq d_i]^{\prod_i S_{d_i}},
 \qquad A=\bigoplus_d A_d.
\]
The charge sum is algebraic. Neither roots nor charges are completed.

Write
\[
 \chi(d,e)=\sum_i d_i e_i-\sum_{a:i\to j}d_i e_j.
\]
In the subobject-first convention, the weighted Hall kernel is
\begin{equation}\label{cyspec:eq:hall-kernel}
 K_{d,e}(x;y)=
 \frac{\displaystyle\prod_{a:i\to j}
                \prod_{r\leq d_i,\ s\leq e_j}(y_{j,s}-x_{i,r}+h_a)}
      {\displaystyle\prod_i
                \prod_{r\leq d_i,\ s\leq e_i}(y_{i,s}-x_{i,r})}.
\end{equation}
The product $F\star G$ sums $F(x_I)G(x_J)K_{d,e}(x_I;x_J)$
over root partitions $I_i\sqcup J_i$ of sizes $d_i,e_i$ at each
vertex. There is no factorial divisor.

This formula is polynomial. A possible pole along a root diagonal
at one vertex is simple. Pair the summands obtained by exchanging
the two roots across the blocks. Their remaining factors agree on
the diagonal, and their residues have opposite signs. Every possible
prime denominator therefore cancels. The result is symmetric.

Both bracketings of three inputs sum over the same ordered
three-block partitions, with one kernel factor for every ordered
pair of blocks. This proves associativity in every charge and,
by the same partition argument, coherence of all iterated products.
The empty block gives the unit $1\in A_0=R$.

Geometrically, the numerator is the normal Euler class of the
invariant-subspace condition. The denominator is the tangent
Euler class of the flag \cite{cyspec:KS}*{\S\S2.4--2.5}.
The product has ordinary cohomological degree $-2\chi(d,e)$.

Define the Euler series, first as a rational function, by
\begin{equation}\label{cyspec:eq:spectral-euler}
 E_{d,e}(u;x,y)=
 \frac{\displaystyle\prod_i\prod_{r\leq d_i,\ s\leq e_i}
                                  (u+y_{i,s}-x_{i,r})}
      {\displaystyle\prod_{a:i\to j}\prod_{r\leq d_i,\ s\leq e_j}
                                  (u+y_{j,s}-x_{i,r}+h_a)}.
\end{equation}
Its leading term is $u^{\chi(d,e)}$.
At $u=0$ it is the inverse of \eqref{cyspec:eq:hall-kernel} in the
ring where the displayed Euler factors are inverted.

Now set
\begin{equation}\label{cyspec:eq:coproduct}
 \Delta_{d,e}(u)F=E_{d,e}(u;x,y)F(x-u,y),
 \qquad F\in A_{d+e},
\end{equation}
where every first-block root is shifted by $-u$.
Expanding at infinity gives a map
$A_{d+e}\to(A_d\otimes_R A_e)((u^{-1}))$.
Each input has finite upper Laurent order, and a fixed total charge
has only finitely many splits. With $\deg u=2$, this component has
cohomological degree $2\chi(d,e)$.

\begin{theorem}\label{cyspec:thm:coassociativity}
On rational representatives the coproduct satisfies
\begin{equation}\label{cyspec:eq:coassociativity}
 (\Delta_{d,e}(u)\otimes1)\Delta_{d+e,f}(v)
 =(1\otimes\Delta_{e,f}(v))\Delta_{d,e+f}(u+v).
\end{equation}
The equality also holds in either chosen iterated Laurent chamber,
using that chamber for every factor on both sides.
\end{theorem}

\begin{proof}
Splitting a block factors its Euler product into the products for
its two constituents. Both sides of \eqref{cyspec:eq:coassociativity}
therefore send $F$ to
\[
 E_{d,e}(u;x,y)E_{d,f}(u+v;x,z)E_{e,f}(v;y,z)
 F(x-u-v,y-v,z).
\]
This is an identity of rational functions before expansion.
For example, one may expand in
$(A_d\otimes_R A_e\otimes_R A_f)((u^{-1}))((v^{-1}))$,
with $v$ the outer variable and $u+v$ expanded at $v=\infty$.
Applying this expansion to the common rational expression proves
the asserted formal identity. It does not require extending either
binary map to every element of an arbitrary two-variable completion.
\end{proof}

Projection onto $A_0$ is the counit $\varepsilon:A\to R$.
The empty-block formulas are
\begin{equation}\label{cyspec:eq:counits}
 (\varepsilon\otimes1)\Delta(u)=\operatorname{id},\qquad
 (1\otimes\varepsilon)\Delta(u)=e^{uD},\qquad
 D=-\sum_{i,r}\partial_{x_{i,r}}.
\end{equation}
The exponential is finite on each polynomial. The even derivation
$D$ lowers cohomological degree by two and satisfies
\[
 \Delta(u)D=\partial_u\Delta(u)+(1\otimes D)\Delta(u).
\]
Indeed, differentiating the Euler factor with respect to $u$ is
the same as differentiating simultaneously in all second-block
roots. These terms cancel against their contribution from $D$.

\subsection{Exchange and the four blocks}

Define
\begin{equation}\label{cyspec:eq:exchange}
 R_{d,e}(u;x,y)=\frac{E_{d,e}(u;x,y)}{E_{e,d}(-u;y,x)}.
\end{equation}
This Laurent unit obeys
$R_{d,e}(u;x,y)R_{e,d}(-u;y,x)=1$.
Its pairwise multipliers on three labeled blocks commute, hence
satisfy the spectral Yang--Baxter identity. Splitting a block gives
the corresponding fusion identities by factorization of $E$.
The signs from negative arguments are included in
\eqref{cyspec:eq:exchange}.

Let $\tau$ exchange two blocks and let $\mathsf T_u$ shift all roots
of a rational representative by $-u$. Then
\begin{equation}\label{cyspec:eq:translated-exchange}
 \Delta_{d,e}(u)=R_{d,e}(u)\mathsf T_u\tau\Delta_{e,d}(-u).
\end{equation}
Simultaneous translation preserves every root difference.
The right side has polynomial part $F(x-u,y)$ and Euler factor
$E_{d,e}(u)$, proving the identity on its stated rational domain.

\begin{theorem}[The four-block identity]\label{cyspec:thm:four-block}
For $F\in A_\alpha$, $G\in A_\beta$ and $r+s=\alpha+\beta$,
\begin{equation}\label{cyspec:eq:four-block}
\begin{split}
 \Delta_{r,s}(u)(F\star G)
 =\sum_{\substack{a+c=\alpha,\ b+d=\beta\\a+b=r,\ c+d=s}}
 &(\star_{a,b}\otimes\star_{c,d})\Bigl[
 R_{b,c}(u;x_b,y_c)\\[-2pt]
 &\quad\cdot\tau_{23}
       (\Delta_{a,c}(u)F\otimes\Delta_{b,d}(u)G)\Bigr].
\end{split}
\end{equation}
Here $\tau_{23}$ takes slots $(a,c,b,d)$ to $(a,b,c,d)$.
After that permutation, $R_{b,c}$ multiplies the middle two slots.
The two Hall products sum over divisions of the output roots.
All operations are formed rationally and then expanded at infinity.
\end{theorem}

\begin{proof}
Write the first-output roots as $x_a,x_b$ and the second-output
roots as $y_c,y_d$. The coproduct on the left contributes
$E_{a,c}E_{a,d}E_{b,c}E_{b,d}$ at parameter $u$.
The Hall kernel between its two input blocks contributes
\[
 K_{a,b}(x_a;x_b),\quad E_{a,d}(u;x_a,y_d)^{-1},\quad
 E_{c,b}(-u;y_c,x_b)^{-1},\quad K_{c,d}(y_c;y_d).
\]

The $a,d$ factors cancel. The $b,c$ factors give
$E_{b,c}(u)/E_{c,b}(-u)=R_{b,c}(u)$.
The remaining Euler factors are those of the two input coproducts.
The remaining kernels are those of the two output products.
Both constructions enumerate the same four labeled subsets.
Thus their summands agree under this correspondence.
\end{proof}

For one vertex and two charge-one inputs, write $K(v)$ for the pair
kernel and put $v=u+y-x$. The constant-input case is already
\[
 \Delta_{1,1}(u)(\mathbf1_1\star\mathbf1_1)
 =\frac{K(v)+K(-v)}{K(v)}=1+R_{1,1}(u;x,y).
\]
The crossed factor is therefore part of the multiplication
comparison itself. The Yang--Baxter equation alone would not
determine this comparison. Equations~\eqref{cyspec:eq:four-block}
and \eqref{cyspec:eq:counits} use the unshifted even coefficient
modules. A different grading must transport their tensor shifts
and permutation signs together.

\section{Critical coefficients}
\label{cyspec:sec:critical}

The same translation mechanism applies when ordinary coefficients
are replaced by vanishing cycles. Let $W$ be a potential, a finite
linear combination of cyclic paths in $Q$, and assume that the
arrow action of $T$ preserves $W$. The arrow weights belong to the
character lattice of $T$. They need not lie in a positive cone.
For $G_d=\prod_iGL_{d_i}$, put
\[
 \mathfrak M_d^T=[V_d/(G_d\times T)],\qquad
 w_d=\operatorname{Tr}(W),\qquad
 B_d=H^\bullet(\mathfrak M_d^T,
                  \phi_{w_d}\mathbb Q[-\chi(d,d)]).
\]

Use the perverse normalization of vanishing cycles and the shift
convention $K[s]^n=K^{n+s}$. The shift $-\chi(d,d)$ is the dimension
of $[V_d/G_d]$, before passage to the external torus quotient.
The coefficient ring remains $R=H^\bullet(BT;\mathbb Q)$.
The charge-dependent tensor convention is
\[
 B_d\otimes^T B_e=(B_d\otimes_R B_e)[\chi(d,e)-\chi(e,d)].
\]

\begin{assumption}[Relative Kunneth]\label{cyspec:ass:kunneth}
For every pair $d,e$, the relative Kunneth map and Thom--Sebastiani
comparison identify
\begin{equation}\label{cyspec:eq:kunneth}
 B_d\otimes_R B_e\xrightarrow{\ \sim\ }
 H^\bullet\left(\mathfrak M_d^T\times_{BT}\mathfrak M_e^T,
 \phi_{w_d\boxplus w_e}\mathbb Q[-\chi(d,d)-\chi(e,e)]\right).
\end{equation}
\end{assumption}

This is an assumption over the base $BT$, imposed component by
component. A Kunneth isomorphism for sheaves on $BT$ does not by
itself give \eqref{cyspec:eq:kunneth} for their cohomology modules.
All the critical coproducts below use this hypothesis.

On the product stack let $E_i$ and $F_i$ be the tautological bundles
of the first and second representation. If $L_a$ is the line on
$BT$ associated with the arrow weight, define the perfect complex
\begin{equation}\label{cyspec:eq:ambient-ext}
 \mathcal V_{d,e}=
 \left[\bigoplus_i E_i^\vee\otimes F_i
 \longrightarrow
 \bigoplus_{a:i\to j}E_i^\vee\otimes F_j\otimes L_a\right],
 \qquad\text{in degrees }0,1.
\end{equation}
The differential sends a family of maps to its commutators with
the arrow maps. This is the derived Hom complex for the underlying
quiver representations. Its virtual rank is $\chi(d,e)$.
It is not the full self-Ext complex in the associated
three-dimensional Calabi--Yau category.

For a virtual bundle $V$, write
\begin{equation}\label{cyspec:eq:characteristic-series}
 \Psi(V,z)=z^{\operatorname{rk}V}c(V,z^{-1}),
 \qquad c(V,t)=\sum_{k\geq0}c_k(V)t^k.
\end{equation}
In a two-term presentation $V=[V^0\to V^1]$, this is the ratio
$\Psi(V^0,z)/\Psi(V^1,z)$. In Chern roots it is the product of
$z+$root for $V^0$, divided by the corresponding product for $V^1$,
expanded at infinity. Thus
\[
 \Psi(V\oplus V',z)=\Psi(V,z)\Psi(V',z),\qquad
 \Psi(V\otimes L,z)=\Psi(V,z+c_1(L)).
\]

Tensoring a representation by a line defines an action of
$B\mathbb G_m$. On coefficients its coaction is
$\operatorname{act}_z^*=e^{zT}$, where $z$ is the first Chern class
of the line and $T$ has degree $-2$.
It sends each Chern root $x$ to $x+z$.
Write $\oplus^*$ for the critical pullback along direct sum,
followed by \eqref{cyspec:eq:kunneth}. The potential identity
$\operatorname{Tr}W(E\oplus F)=w_d(E)+w_e(F)$ supplies the
Thom--Sebastiani comparison used in this map.

A \emph{coassociative vertex coproduct} here means a Laurent map
satisfying parameter-dependent coassociativity
\eqref{cyspec:eq:coassociativity} and the translation identity.
It also satisfies the two counit identities \eqref{cyspec:eq:counits}.
Use $z,T$ in place of $u,D$ and the specified tensor convention.
Its right counit is polynomial on each input.
This definition imposes no exchange axiom.

\begin{theorem}[The critical spectral coproduct]
\label{cyspec:thm:critical-coproduct}
Under Assumption~\ref{cyspec:ass:kunneth}, the maps
\begin{equation}\label{cyspec:eq:critical-coproduct}
 \Delta^{\mathrm{crit}}_{d,e}(z)
 =\Psi(\mathcal V_{d,e},-z)\operatorname{act}_{z,1}^*\oplus^*
\end{equation}
give a coassociative vertex coproduct on $B=\bigoplus_d B_d$.
Its component has target
\[
 (B_d\otimes_R B_e)[\chi(d,e)-\chi(e,d)]((z^{-1})).
\]
The covacuum is projection to $B_0=R$. Its left counit is the
identity and its right counit is $e^{zT}$.
\end{theorem}

\begin{proof}
The direct-sum pullbacks are coassociative and commute with
simultaneous line tensoring. Pulling $\mathcal V$ back along a
direct sum in either variable splits it into the corresponding
two Ext complexes. Tensoring the first variable by the line changes
its characteristic argument by $-z$.
Consequently the two three-block composites have common factor
\[
 \Psi(\mathcal V_{d,e},-z)
 \Psi(\mathcal V_{d,f},-z-w)
 \Psi(\mathcal V_{e,f},-w).
\]
They tensor the first block by the line of class $z+w$
and the second by the line of class $w$.

This proves the coassociativity identity with parameters $z,w,z+w$,
in compatible Laurent expansions. The empty Ext complex gives the
two counit identities. Differentiating the line action gives
$\Delta^{\mathrm{crit}}(z)T=\partial_z\Delta^{\mathrm{crit}}(z)
+(1\otimes T)\Delta^{\mathrm{crit}}(z)$.

The perverse shifts give direct-sum pullback degree
$-\chi(d,e)-\chi(e,d)$. Multiplication by the characteristic factor
adds $2\chi(d,e)$. The stated target shift makes the resulting map
degree zero. These are the geometric and grading identities of
\cite{cyspec:JKL}*{Theorem 2.3.1}.
\end{proof}

The localized operation before line translation is
\[
 \Delta^{\mathrm{loc}}_{d,e}
 =e(\mathcal V_{d,e})\oplus^*,\qquad
 e(\mathcal V_{d,e})=
 \frac{e(\mathcal V_{d,e}^0)}{e(\mathcal V_{d,e}^1)}.
\]
Inverting the indicated Euler classes makes this expression defined.
Applying $\operatorname{act}_{z,1}^*$ and expanding at infinity
changes $e(\mathcal V)$ to $\Psi(\mathcal V,-z)$.
This identifies the localized and spectral constructions
\cite{cyspec:JKL}*{Lemma 3.2.4 and Theorem 3.3.2}.
For zero potential, the unshifted polynomial formula
\eqref{cyspec:eq:coproduct} uses $z=-u$.

\subsection{Braided multiplication and its signs}

Assume now that $Q$ is symmetric: for each pair of distinct vertices
the numbers of arrows in the two directions agree. Thus $\chi$ is
symmetric, although the arrow weights need not be paired by negation.
The tensor shift in Theorem~\ref{cyspec:thm:critical-coproduct} vanishes.
Let $H$ be the tautological ring, the $R$-subalgebra of
$\prod_d H^\bullet(\mathfrak M_d^T)$ generated by the Chern characters
of the vertex bundles, including their ranks.
Direct sum and line tensoring give its coproduct
$\delta_H(z)=\operatorname{act}_{z,1}^*\oplus^*$.
It acts on $B$ by cup product.

Let $\sigma$ denote the graded swap, with its Koszul sign on
cohomology. The characteristic multiplier
\begin{equation}\label{cyspec:eq:joyce-r}
 \mathscr R_{d,e}(z)=
 \frac{\Psi(\sigma^*\mathcal V_{e,d}^{\vee},z)}
      {\Psi(\mathcal V_{d,e},z)}
\end{equation}
has coefficients in $H\otimes_R H$ and leading term one.
Its spectral hexagon identities follow from direct-sum factorization
and line translation of the two complexes.
These identities define the meromorphic exchange
$\beta(z)=\sigma\circ(\mathscr R(z)\cdot-)$ on tautological modules,
using $\delta_H(z)$ for the coefficient action.

In Chern-root coordinates its relation to
\eqref{cyspec:eq:exchange} is
\begin{equation}\label{cyspec:eq:exchange-conventions}
 \mathscr R_{e,d}(-u;y,x)=(-1)^{\chi(d,e)}R_{d,e}(u;x,y).
\end{equation}
Indeed, duality gives $\Psi(V^\vee,z)=(-1)^{\operatorname{rk}V}
\Psi(V,-z)$. Thus the change $z=-u$ includes both a block reversal
and the displayed Euler sign.

The perverse grading contributes another sign. Put
\begin{equation}\label{cyspec:eq:tau}
 p(d)=\chi(d,d)\pmod2,\qquad
 \tau(d,e)=p(d)p(e)+\chi(d,e)\pmod2.
\end{equation}
Symmetry of $\chi$ makes $p$ linear modulo two and makes $\tau$
alternating and bilinear. Use the exchange
\[
 \beta^\tau_{d,e}(z)=(-1)^{\tau(d,e)}\beta_{d,e}(z).
\]

The critical Hall product $m$ and the coproduct satisfy
\begin{equation}\label{cyspec:eq:critical-bialgebra}
 \Delta^{\mathrm{crit}}(z)m
 =(m\otimes m)(1\otimes\beta^\tau(z)\otimes1)
        (\Delta^{\mathrm{crit}}(z)\otimes\Delta^{\mathrm{crit}}(z)).
\end{equation}
This is the untwisted-product form of the critical vertex-bialgebra
theorem \cite{cyspec:JKL}*{Theorem 4.5.2 and its untwisted form}.
Its proof starts with the localized Hall compatibility, translates
the first and third blocks, and expands their invertible Euler
factors. Cup-product linearity permits translation through the two
Hall products. The crossed Euler factor becomes
\eqref{cyspec:eq:joyce-r}. The Koszul and perverse signs give
\eqref{cyspec:eq:tau}.

One can move this sign into the operations, provided both are
transported. Order the vertices, with basis dimension vectors $e_i$,
and choose
\[
 \psi(d,e)=\sum_{i<j}\tau(e_i,e_j)d_i e_j\pmod2.
\]
Then $\psi(d,e)+\psi(e,d)=\tau(d,e)$ and $\psi$ is bilinear.

\begin{proposition}[Simultaneous sign transport]
\label{cyspec:prop:sign-transport}
Define, on their charge components,
\begin{equation}\label{cyspec:eq:sign-transport}
 m^\psi_{d,e}=(-1)^{\psi(d,e)}m_{d,e},\qquad
 \Delta^\psi_{d,e}(z)=(-1)^{\psi(d,e)}
                         \Delta^{\mathrm{crit}}_{d,e}(z).
\end{equation}
They remain unital, associative and spectrally coassociative,
with the same counits and translation. Their multiplication
compatibility uses $\beta$ in place of $\beta^\tau$.
\end{proposition}

\begin{proof}
Bilinearity gives
$\psi(d,e)+\psi(d+e,f)=\psi(e,f)+\psi(d,e+f)$,
so the two associativity identities retain their signs.
Also $\psi(d,0)=\psi(0,d)=0$.

For four blocks set $r=a+b$, $s=c+d$, $\alpha=a+c$, $\beta=b+d$.
Expanding the bilinear terms gives
\[
 \psi(r,s)+\psi(\alpha,\beta)+\tau(c,b)
 =\psi(a,b)+\psi(c,d)+\psi(a,c)+\psi(b,d)\pmod2.
\]
The left side is the sign on the left of
\eqref{cyspec:eq:critical-bialgebra}, including its old crossed sign.
The right side is the sign of the four transported operations.
This proves the assertion.
\end{proof}

Changing only multiplication need not preserve this square.
For two isolated vertices, $\chi$ is the identity matrix and
$\psi(e_1,e_2)=1$. The straight four-block term
$(a,c,b,d)=(e_1,0,0,e_2)$ changes sign on the product of the two
inputs, while neither product in the two outputs changes sign.
The coproduct factors in \eqref{cyspec:eq:sign-transport} supply
the required correction.

For a nonsymmetric quiver, the characteristic ratio
\eqref{cyspec:eq:joyce-r} has leading power
$z^{\chi(e,d)-\chi(d,e)}$. The componentwise module formulation
of \cite{cyspec:JKL}*{\S4.5.4} therefore has a different coefficient
recipient from the tautological-ring formulation above.

The displayed symmetric sign prescription also has a precise limit.
For one arrow $1\to2$,
\[
 \chi(e_1,e_2)=-1,\quad\chi(e_2,e_1)=0,\quad
 \tau(e_1,e_2)=0,\quad\tau(e_2,e_1)=1.
\]
No function $\psi$ can then satisfy $\psi+\psi^{\mathrm{op}}=\tau$.
Extending that sign convention requires a compatible choice on the
shifted componentwise tensor product. The polynomial identity
\eqref{cyspec:eq:four-block}, with its explicitly even convention,
already applies to this quiver.

\section{An exchange ratio and two integral lattices}
\label{cyspec:sec:lattices}

Exchange compares a pair kernel with its reverse. Multiplying the
kernel by an even function leaves that ratio unchanged, but can
change multiplication on the collision divisor. To measure the
change, keep the entire polynomial coefficient module in each charge.

For this section take
\begin{equation}\label{cyspec:eq:parameter-ring}
 R=\mathbb Q[h_1,h_2],\qquad h_3=-h_1-h_2,\qquad
 a=\sum_{i<j}h_i h_j,\qquad b=h_1h_2h_3.
\end{equation}
Thus $a=-(h_1^2+h_1h_2+h_2^2)$ and $b=-h_1h_2(h_1+h_2)\ne0$
in the domain $R$. No parameter is inverted.

Define two independently specified polynomial carriers
\[
 A_n^{(j)}=R[x_1,\ldots,x_n]^{S_n},\qquad
 \mathcal A^{(j)}=\bigoplus_{n\geq0}A_n^{(j)},\qquad j=0,1,
\]
with even parity and algebraic charge sum. Their pair kernels are
\begin{equation}\label{cyspec:eq:two-kernels}
 K_0(v)=\frac{\prod_{i=1}^3(v+h_i)}{v}
       =v^2+a+\frac bv,\qquad
 K_1(v)=v^2K_0(v).
\end{equation}
Here $v$ is a root difference, not a coefficient parameter.

For $F\in A_d^{(j)}$ and $G\in A_e^{(j)}$, set
\begin{equation}\label{cyspec:eq:two-products}
 (F*_jG)(X)=\sum_{\substack{I\sqcup J=[d+e]\\|I|=d}}
 F(X_I)G(X_J)\prod_{r\in I,\ s\in J}K_j(x_s-x_r).
\end{equation}
The sum has no factorial divisor and the empty charge is its unit.

The pole-cancellation and ordered-partition arguments of
Section~\ref{cyspec:sec:polynomial-coproduct} prove polynomiality and
associativity. For $K_1$, polynomiality also follows term by term.
The first kernel is the zero-potential three-loop Hall kernel.
Two additional loops of weight zero give the second kernel.
This is a calculation of ordinary normal Euler classes, without
a critical-cohomology identification.

Both kernels have exchange ratio
\begin{equation}\label{cyspec:eq:common-ratio}
 \frac{K_0(-v)}{K_0(v)}=\frac{K_1(-v)}{K_1(v)}
 =\frac{v^3+av-b}{v^3+av+b}
 =\frac{\prod_i(v-h_i)}{\prod_i(v+h_i)}.
\end{equation}
The equality cancels the even factor $v^2$.
To decide whether that cancellation comes from an invertible
change of polynomial coefficients, compute the whole binary image.

\subsection{Binary generation and its cokernel}

In charge two put
\[
 c=\frac{x+y}{2},\qquad v=y-x,\qquad t=v^2,\qquad B=R[c,t].
\]
Then $A_2^{(j)}=B$ and
$A_1^{(j)}\otimes_R A_1^{(j)}=R[c,v]=B\oplus vB$.
Pure tensors span this entire source. For $p=A+vC$, with $A,C\in B$,
linearized multiplication is therefore
\begin{equation}\label{cyspec:eq:binary-rows}
 \mu_0(A+vC)=2((t+a)A+bC),\qquad
 \mu_1(A+vC)=2t((t+a)A+bC).
\end{equation}

Indeed, the two terms for $K_0$ are
$(v^2+a+b/v)(A+vC)$ and $(v^2+a-b/v)(A-vC)$.
Their sum gives the first formula. The common even factor $v^2$
gives the second. Since $2$ is invertible, their images are
\begin{equation}\label{cyspec:eq:binary-images}
 \operatorname{im}\mu_0=J=(t+a,b)B,\qquad
 \operatorname{im}\mu_1=tJ.
\end{equation}

Regard each product as the two-term cochain complex
\[
 M_j=[B\oplus vB\xrightarrow{\mu_j}B],
 \qquad\text{in degrees }-1,0.
\]
It records binary generation, rather than a full deformation complex.

\begin{proposition}\label{cyspec:prop:binary-cohomology}
The cohomology of these complexes is
\begin{align}
 H^{-1}(M_0)=H^{-1}(M_1)&=B\,(b,-(t+a)),
       \label{cyspec:eq:binary-kernel}\\
 H^0(M_0)&=R[c]/(b),\label{cyspec:eq:first-cokernel}\\
 H^0(M_1)&=R[c]\oplus (R[c]/(b))\,t
       \quad\text{as }R[c]\text{-modules}.
       \label{cyspec:eq:second-cokernel}
\end{align}
On the last module, multiplication by $t$ is
$t(A,C)=(0,A-aC)$, with the second coordinate reduced modulo $b$.
Evaluation at $t=0$ gives a cochain map $M_1\to R[c][0]$.
It detects the nonzero class of $b\mathbf1_2$.
\end{proposition}

\begin{proof}
If $(t+a)A+bC=0$, reduction modulo $t+a$ gives $b\overline C=0$
in the domain $R[c]$. Thus $C=-(t+a)H$ and $A=bH$.
Multiplication of the row by the non-zero-divisor $t$ changes no
kernel, proving \eqref{cyspec:eq:binary-kernel}.

The first cokernel sets $t=-a$ and $b=0$.
In the second, the relations are $t^2=-at$ and $bt=0$.
Division by the monic polynomial $t^2+at$ gives a unique expression
$A(c)+tC(c)$, with $C$ then taken modulo $b$.
This proves both the module description and its $t$-action.

Every value of $\mu_1$ is divisible by $t$, so evaluation at zero
annihilates the differential. It sends $b\mathbf1_2$ to the nonzero
element $b$. By contrast, $\mu_0(v/2)=b\mathbf1_2$.
\end{proof}

An \emph{ambient lattice equivalence} will mean a family of
$R$-linear bijections $U_n:A_n^{(0)}\to A_n^{(1)}$ with $U_0=\operatorname{id}_R$
and
\begin{equation}\label{cyspec:eq:ambient-gauge}
 U_{d+e}(F*_0G)=U_d(F)*_1U_e(G).
\end{equation}
No preservation of root degree, translation or multiplication by
$c,t$ is assumed. In particular, these are not restricted to maps
given by scalar multiplication.

\begin{theorem}\label{cyspec:thm:ambient-obstruction}
There is no ambient lattice equivalence. The conclusion remains
true if every $U_n$ is semilinear for one automorphism of $R$.
\end{theorem}

\begin{proof}
The $R$-annihilators of the binary cokernels are
\[
 \operatorname{Ann}_R H^0(M_0)=(b),\qquad
 \operatorname{Ann}_R H^0(M_1)=0.
\]
The first follows from \eqref{cyspec:eq:first-cokernel}. The second
follows from the quotient $H^0(M_1)\twoheadrightarrow R[c]$, a
faithful $R$-module. The charge-two instance of
\eqref{cyspec:eq:ambient-gauge} would identify the complexes through
$U_1\otimes U_1$ in degree $-1$ and $U_2$ in degree zero.
Their cokernels would then have the same annihilator, a contradiction.

Under a parameter automorphism $\sigma$, the nonzero annihilator
$b$ is carried to the nonzero element $\sigma(b)$. The contradiction
persists. No step requires $U_2$ to preserve the coordinate $t$.
\end{proof}

The same distinction survives inversion of parameters alone.
For a submodule $L\subset B$ define its parameter saturation by
\[
 L_R^{\mathrm{sat}}=\{F\in B:rF\in L
                     \text{ for some }0\ne r\in R\}.
\]
Since $bB\subset J$ and $btB\subset tJ$,
\begin{equation}\label{cyspec:eq:saturation}
 J_R^{\mathrm{sat}}=B,\qquad (tJ)_R^{\mathrm{sat}}=tB.
\end{equation}

For the reverse inclusion in the second equality, if $rF\in tB$,
evaluation at $t=0$ gives $rF(c,0)=0$ in the domain $R[c]$.
Thus $F\in tB$. Over $\operatorname{Frac}(R)$, the first binary
product is consequently surjective and the second still has image
$t\operatorname{Frac}(R)[c,t]$. In particular,
$\mu_0(v/(2b))=\mathbf1_2$ but $\mu_1(v/(2b))=t\mathbf1_2$.

\subsection{The comparison that does exist}

Set $D_0=D_1=1$ and
\begin{equation}\label{cyspec:eq:discriminant}
 D_n(X)=\prod_{r<s}(x_s-x_r)^2,\qquad \Phi_n(F)=D_nF.
\end{equation}
For each partition $I\sqcup J=[d+e]$,
\begin{equation}\label{cyspec:eq:discriminant-factorization}
 D_{d+e}(X)=D_d(X_I)D_e(X_J)
                  \prod_{r\in I,\ s\in J}(x_s-x_r)^2.
\end{equation}
Thus the even factors omitted by exchange supply an actual map.

\begin{theorem}\label{cyspec:thm:discriminant-map}
The maps $\Phi_n$ define a unital, injective, charge-preserving
algebra map $\Phi:\mathcal A^{(0)}\to\mathcal A^{(1)}$.
They commute with simultaneous root translation and satisfy
\begin{equation}\label{cyspec:eq:discriminant-product}
 \Phi_{d+e}(F*_0G)=\Phi_d(F)*_1\Phi_e(G).
\end{equation}
They are not surjective on $A_n$ for $n\geq2$.
If $\mathcal C^{(j)}$ is the subalgebra generated by $A_1^{(j)}$,
restriction gives an integral isomorphism
\begin{equation}\label{cyspec:eq:current-isomorphism}
 \Phi:\mathcal C^{(0)}\xrightarrow{\sim}\mathcal C^{(1)},
 \qquad\mathcal C_n^{(1)}=D_n\mathcal C_n^{(0)}.
\end{equation}
\end{theorem}

\begin{proof}
In every summand on the right of
\eqref{cyspec:eq:discriminant-product}, the internal discriminants
and the cross squares from $K_1=v^2K_0$ combine, by
\eqref{cyspec:eq:discriminant-factorization}, into $D_{d+e}$.
This factor is independent of the subset and can be taken outside
the sum. The same factorization over any ordered partition proves
compatibility with every iterated product.

Each $D_n$ is nonzero in a domain, so multiplication by it is
injective. For $n\geq2$ its image cannot contain the constant
polynomial, since $D_n$ vanishes at $x_1=x_2$.
The factors $D_0=D_1=1$ give the unit and the identity on charge
one. A multiplicative injection fixing those generators therefore
maps $\mathcal C^{(0)}$ exactly onto $\mathcal C^{(1)}$.
Division by $D_n$ on this specified image gives its inverse.
Finally $\sum_i\partial_{x_i}D_n=0$, proving translation compatibility.
\end{proof}

This isomorphism of generated algebras and
Theorem~\ref{cyspec:thm:ambient-obstruction} concern different modules.
The former identifies the generated images, whereas the latter
excludes its extension to bijections of all polynomial coefficients.
For compatible even gradings, give roots and parameters raw degree
two and set
\[
 \deg_0F=\deg_{\mathrm{raw}}F-2n(n-1),\qquad
 \deg_1F=\deg_{\mathrm{raw}}F-4n(n-1),\qquad F\in A_n.
\]
Both products and $\Phi$ then have degree zero. The obstruction
requires none of these grading restrictions on the proposed gauges.

Define the spectral maps of the two algebras independently by
\begin{equation}\label{cyspec:eq:lattice-coproduct}
 \Delta^{(j)}_{d,e}(s)F(X,Y)=
 \prod_{r,k}K_j(s+y_k-x_r)^{-1}\,F(X-s,Y).
\end{equation}
They are rationally coassociative by
Theorem~\ref{cyspec:thm:coassociativity}, and expand at infinity with
coefficients in $A_d^{(j)}\otimes_R A_e^{(j)}$.
Their left counit is the identity and their right counit is
translation by $-s$.

\begin{proposition}\label{cyspec:prop:localized-comparison}
The discriminant maps preserve these coproducts:
\begin{equation}\label{cyspec:eq:discriminant-coproduct}
 \Delta^{(1)}_{d,e}(s)\Phi_{d+e}
 =(\Phi_d\otimes_R\Phi_e)\Delta^{(0)}_{d,e}(s).
\end{equation}
The products extend to
$\mathcal A^{(j)}_{\mathrm{reg}}=\bigoplus_n A_n^{(j)}[D_n^{-1}]$.
On these carriers $\Phi$ is an algebra isomorphism preserving
translation and the coproducts with localized output coefficients.
\end{proposition}

\begin{proof}
The translated factorization is
\[
 D_{d+e}(X-s,Y)=D_d(X)D_e(Y)
                         \prod_{r,k}(s+y_k-x_r)^2.
\]
The cross factors cancel the extra squares in the denominators of
$\Delta^{(1)}$, proving \eqref{cyspec:eq:discriminant-coproduct}
before and after expansion.

Every internal discriminant and every cross-root difference is
invertible after the total discriminant is inverted. Each localized
shuffle therefore has poles only on that total divisor and lies in
the stated localized symmetric ring. Multiplication by $D_n$ now
has inverse $D_n^{-1}$. Translation fixes $D_n$, and the displayed
translated factorization proves the coproduct assertion.
In charge two the map of binary complexes $(\operatorname{id},t)$
is an isomorphism after $t$ is inverted. No coefficient parameter
such as $b$ needs to be inverted.
\end{proof}

\subsection{The remaining collision module}

The classes discarded by the discriminant localization form the
coefficient modules
\[
 Q_n=A_n^{(1)}/D_nA_n^{(1)},\qquad Q=\bigoplus_n Q_n.
\]
Their support lies on the internal root discriminant.
In particular, $Q_0=Q_1=0$ and $Q_2=B/(t)=R[c]$.
There is an exact sequence
\begin{equation}\label{cyspec:eq:relative-sequence}
 0\longrightarrow\mathcal A^{(0)}\xrightarrow{\Phi}\mathcal A^{(1)}
 \longrightarrow Q\longrightarrow0
\end{equation}
of $\mathcal A^{(0)}$-bimodules, with quotient actions
\begin{equation}\label{cyspec:eq:quotient-actions}
 a\cdot[z]=[\Phi(a)*_1z],\qquad [z]\cdot a=[z*_1\Phi(a)].
\end{equation}

These are well-defined because replacing $z$ by $z+\Phi(c)$ changes
the two products by $\Phi(a*_0c)$ and $\Phi(c*_0a)$.
Associativity proves the bimodule laws. Translation preserves the
sequence, since it fixes $D_n$ and is a derivation of both products.
The sequence is not an algebra quotient: its proper kernel contains
the unit and therefore cannot be an ideal of a nonzero unital quotient.

The charge-two cone is $[B\xrightarrow{t}B]$ in degrees $-1,0$,
with cokernel $B/(t)$. The same cone occurs for the chain map
$(\operatorname{id},t):M_0\to M_1$.
Indeed, write $P=B^2$. Its cone has terms $P$, $P\oplus B$, $B$
in degrees $-2,-1,0$, with differentials
\[
 p\longmapsto(p,-\mu_0p),\qquad
 (q,z)\longmapsto t(\mu_0q+z).
\]
The coordinate $w=z+\mu_0q$ splits off the contractible pair
$P\xrightarrow{1}P$, leaving $[B\xrightarrow{t}B]$.

On the ordered root cover this quotient is $R[c,v]/(v^2)$.
Its normal direction has length two there, while on the symmetric
quotient its coordinate $t=v^2$ gives length one.

\subsection{Other even factors}

The same argument applies to a general monic even multiplier.
Let $R$ now be a commutative $\mathbb Q$-algebra that is a domain,
and choose $a,b\in R$ with $b\ne0$.
Put $q(v)=P(v^2)$, where $P(t)\in R[t]$ is monic of positive degree.

\begin{theorem}\label{cyspec:thm:general-even-factor}
The shuffle products for $K(v)=v^2+a+b/v$ and $q(v)K(v)$ are
polynomial and associative. There is no charge-preserving
$R$-linear equivalence of their full polynomial algebras.
Their current-generated subalgebras are nevertheless isomorphic
through multiplication in charge $n$ by
\begin{equation}\label{cyspec:eq:general-pair-factor}
 \prod_{r<s}q(x_s-x_r).
\end{equation}
This map is injective on the full carriers and becomes an algebra
isomorphism after the displayed pair product is inverted.
The two exchange ratios agree.
\end{theorem}

\begin{proof}
Pole cancellation over $\operatorname{Frac}(R)$ removes the possible
simple diagonal denominators. Division by each monic linear root
difference already takes place over $R$. Successive such divisions
remove the full denominator: over the fraction field, distinct
differences are relatively prime. Ordered partitions then prove
associativity over $R$.

In charge two the images are $J=(t+a,b)$ and $P(t)J$ in $B=R[c,t]$.
The first cokernel is killed by $b$. The second surjects onto
$B/(P(t))$, which is a nonzero free $R[c]$-module with basis
$1,t,\ldots,t^{\deg P-1}$. It is faithful over $R$.
The scalar-annihilator argument therefore excludes every proposed
equivalence of full carriers. Their parameter saturations are
$B$ and $P(t)B$: use $bP(t)B\subset P(t)J$ and the
$R$-torsion-freeness of $B/P(t)B$.

The pair product \eqref{cyspec:eq:general-pair-factor} factors into
internal and cross-block products. Substitution in each shuffle
summand proves multiplicativity and the isomorphism on generated
subalgebras. It is injective over the domain.

For localized inputs, a sufficiently high power of the total pair
product clears their internal denominators. The remaining diagonal
poles cancel by the polynomiality argument. This proves closure
under localized multiplication and gives the stated inverse.
It does not require that inverting $q(v)$ invert $v$.
Finally $q(-v)=q(v)$ proves equality of the exchange ratios.
\end{proof}

At $(h_1,h_2,h_3)=(1,1,-2)$, the original example has $a=-3$, $b=-2$.
Its two constant-current products are $2(t-3)$ and $2t(t-3)$.
The tensor polynomial $-v/4$ multiplies to $\mathbf1_2$ for $K_0$
and to $t\mathbf1_2$ for $K_1$.

At a fiber where $b=0$, the broad class of linear gauges requires
a separate invariant. For example, when $a=b=0$, the two cokernels
are $\mathbb Q[c,t]/(t)$ and $\mathbb Q[c,t]/(t^2)$.
Their $t$-module structures differ, but an arbitrary
$\mathbb Q$-linear gauge need not preserve that operator.
The integral theorem asserts no such fiberwise classification.

\section{From coefficient collisions to fields}
\label{cyspec:sec:realization}

The coefficient constructions now retain their operations:
direct-sum Gysin maps and their cones, Hall multiplication,
spectral translation, the four-block square, and the discriminant
comparison. The even-factor example shows why a common exchange
ratio cannot recover the coefficient module supported at a collision.
It also supplies a positive comparison on generated algebras and
on the complement of that collision.

\begin{problem}\label{cyspec:prob:field-realization}
Construct a geometric coefficient family on a spacetime collision
space which realizes these spectral operations. Its comparison with
the Chern-root family must preserve Hall multiplication, translation,
the empty-charge unit and the chosen local vacuum. Across the
collision divisor it must realize the derived coefficient cones
and their attachment maps.

For \eqref{cyspec:eq:relative-sequence}, construct both bimodule
actions \eqref{cyspec:eq:quotient-actions} geometrically and identify
the first relative term with $[R[c,t]\xrightarrow{t}R[c,t]]$,
including the quotient class detected by $t=0$.
At three or more points, require compatibility on all intersecting
collision divisors and under every successive merger.
\end{problem}

This problem requires a map between two independently specified
families. Internal root differences have so far been coefficient
coordinates. Their substitution for spacetime separations needs
pullback, support and continuity data. The translated counit of the
spectral coproduct also remains part of that comparison. It does
not become an ordinary Hopf counit by choosing an expansion.
